\chapter{The analysis procedure, and its behaviour on constructed data}
\label{app:analysis}
\label{ch:participant}

\noindent\emph{The analysis \cref{app:study} would be evaluated with,
measured against data whose answer is chosen. No dyad data exists yet; this
is a property of the procedure, not a result.}

The design is within dyads: for task $j$, the paired difference
$d_{ij}=y^{\text{SA}}_{ij}-y^{\text{DT}}_{ij}$ and its standardised mean
$d_z=\bar{d}_{\cdot j}/s_{d,j}$ is the quantity the design was powered for
(distinct from the standardised difference of means, which deflates every
effect by $\sqrt2$ or more once within-dyad correlation is accounted for).
Five tasks give five comparisons, corrected with Holm's step-down procedure,
which controls the family-wise error rate under any dependence structure and
is uniformly more powerful than Bonferroni.

Applied to \num{20000} simulated studies at the planned sixteen dyads (a
true $d_z=0.73$ effect on all five contrasts; and separately, no effect at
all), through the unmodified analysis code:

\begin{table}[htbp]
  \centering
  \caption[The analysis on constructed data]{$n=16$, $m=5$, $\alpha=0.05$.}
  \label{tab:pipeline-validation}
  \begin{tabular}{@{}lccc@{}}
    \toprule
    & one contrast, & one contrast, & some contrast \\
    Constructed truth & uncorrected & Holm-corrected & in the family \\
    \midrule
    $d_z=0.73$ on all five & \num{0.777} & \num{0.638} & \num{0.973} \\
    no effect at all       & \num{0.052} & \num{0.010} & \num{0.050} \\
    \bottomrule
  \end{tabular}
\end{table}

The null row is the check that could fail: a false-positive rate of
\num{0.050} against nominal $\alpha=0.05$ shows the correction is
implemented as specified. The design's own \SI{80}{\percent} power figure is
confirmed \emph{uncorrected} (\SI{77.7}{\percent}) but no contrast will be
reported uncorrected; after Holm, recovering the effect on a \emph{named}
task falls to \SI{63.8}{\percent}, and restoring \SI{80}{\percent} on a named
task needs twenty dyads, not sixteen. The study is nevertheless well powered
for the question it actually asks -- whether shared autonomy changes
performance on this platform at all, where the family-wise power is
\SI{97.3}{\percent} -- so a per-task null at $n=16$ must be reported as an
interval, not an absence. Trial validity is defined in one shared module
that exits non-zero when no trials survive or more than half are excluded;
before that existed, a session with every trial invalidated analysed
cleanly to empty means and undefined intervals, indistinguishable from a
null result on inspection.
